\chapter{Pruebas}
	\label{cap:pruebas}
	Este capítulo documenta la verificación de la plataforma: la estrategia con que se diseñaron las
	pruebas y los casos y resultados de cada característica de calidad. Las pruebas se organizan siguiendo el modelo ISO/IEC 25010 \cite{iso25010} con que se
	clasificaron los requerimientos no funcionales (Tabla~\ref{tab:rnf}), en siete secciones:
	funcionalidad, mantenibilidad, usabilidad, fiabilidad, seguridad, capacidad y confidencialidad. Cada
	caso de prueba indica el requerimiento que verifica, lo que permite rastrear cada requerimiento
	hasta la evidencia que lo respalda.

	Una advertencia de método acompaña a todo el capítulo: probar no demuestra la ausencia de defectos,
	solo reduce el riesgo de que existan \cite{myers_testing}. Por eso las pruebas se priorizaron según
	el daño que causaría cada falla, y los resultados se reportan tal como se obtuvieron, incluidas las
	pruebas que fallaron y las que siguen pendientes de ejecutarse.

	\section{Estrategia de pruebas}
	\label{sec:estrategia_pruebas}
	\subsection{Niveles y técnicas}
	La estrategia sigue el modelo en V, en el que cada nivel de prueba verifica el artefacto del nivel de
	diseño correspondiente \cite{istqb_foundation}:

	\begin{itemize}
		\item \textbf{Pruebas de componente (unitarias).} Verifican una clase de reglas o de servicio de
		forma aislada, sustituyendo sus dependencias por objetos simulados. Se aplican solo a la lógica que
		toma decisiones; no se prueban métodos de acceso, DTO ni repositorios generados.
		\item \textbf{Pruebas de integración.} Verifican la colaboración de un servicio con su base de
		datos, su configuración de seguridad y los servicios que invoca. Se siguió una estrategia
		ascendente: primero cada servicio con su base, después los pares de servicios.
		\item \textbf{Pruebas de sistema.} Recorren los casos de uso de extremo a extremo sobre el sistema
		desplegado en un entorno de pruebas.
		\item \textbf{Pruebas de aceptación.} Ejecutadas por personal de la Defensoría, o por personas que
		actúan cada rol, bajo supervisión del equipo (pruebas alfa).
	\end{itemize}

	Según la información disponible para diseñarlas, las pruebas son de \textbf{caja blanca} ---se
	diseñan a partir del código, como las unitarias y la verificación estática de permisos--- o de
	\textbf{caja negra} ---se diseñan a partir de la especificación y se ejecutan contra la interfaz
	pública, como las de seguridad dinámica y de confidencialidad---. Según su forma de ejecución, son
	\textbf{estáticas}, si analizan el código o la configuración sin ejecutarlos, o
	\textbf{dinámicas}, si ejecutan el sistema. Para diseñar los datos de entrada se emplearon las
	técnicas de particiones de equivalencia, análisis de valores límite, tablas de decisión y transición
	de estados \cite{myers_testing}; por ejemplo, el plazo de 90 días se prueba con 89, 90 y 91 días, y
	la máquina de estados de la queja con cada transición permitida y con las prohibidas.

	\subsection{Prioridad basada en riesgo}
	No es viable probar todas las combinaciones de entrada, de modo que el esfuerzo se asignó según el
	daño que causaría cada falla. Son de prioridad P1 los riesgos de acceso indebido a datos personales,
	de ruptura del ciclo de estatus, de atención simultánea de una queja por dos recepcionistas, de falla
	en el paso de expedientes entre áreas y de exposición del tráfico o de los secretos; de prioridad
	P2, los errores en plazos y datos obligatorios, la pérdida de respaldos y la caída de servicios
	auxiliares; y de prioridad P3, la lentitud y las pantallas confusas.

	\subsection{Entorno y criterios}
	Las pruebas dinámicas se ejecutan en un entorno de pruebas y no en producción: una base PostgreSQL~16
	dedicada, sembrada con los datos de referencia y con la carga masiva de quejas sintéticas, los mismos
	artefactos que se despliegan y una cuenta de prueba por rol. Las pruebas estáticas y unitarias no
	necesitan el sistema levantado. Se definieron los siguientes criterios:

	\begin{itemize}
		\item \textbf{Entrada:} el código compila sin errores y el entorno de pruebas está disponible.
		\item \textbf{Salida:} todos los casos P1 ejecutados y aprobados, al menos el 90~\% de los P2
		ejecutados, ningún defecto crítico o alto abierto, cobertura de al menos 70~\% en las clases de
		reglas, percentil 95 menor a 1~s en la prueba de carga y puntaje SUS de al menos 68.
		\item \textbf{Suspensión:} si un defecto crítico bloquea más del 30~\% de los casos, se detiene la
		ejecución hasta corregirlo.
	\end{itemize}

	\section{Pruebas de funcionalidad}
	\label{sec:pruebas_funcionalidad}
	Las pruebas de funcionalidad verifican la adecuación funcional: que las reglas de negocio y los
	casos de uso se comporten como se especificaron.

	\subsection{Pruebas unitarias}
	Se escribieron pruebas unitarias con JUnit~5 y Mockito para los diez microservicios Java
	desarrollados para la plataforma \cite{junit5, mockito}. Las pruebas no levantan el contexto de
	Spring ni requieren base de datos o red: las dependencias se sustituyen por objetos simulados, lo que
	permite ejecutarlas en segundos en cada compilación. Siguen el patrón preparar--actuar--verificar, incluyen
	un mensaje en cada aserción y se diseñaron con particiones de equivalencia y valores límite. La
	Tabla~\ref{tab:pruebas_unitarias} resume su alcance.

	\rowcolors{2}{rowblue}{white}
	\small
	\renewcommand{\arraystretch}{1.2}
	\setlength{\tabcolsep}{4pt}
	\begin{longtable}{|p{3.4cm}|L{1.3cm}|L{9.8cm}|}
		\hiderowcolors\caption{Pruebas unitarias por microservicio}
		\label{tab:pruebas_unitarias} \\ \showrowcolors
		\hline
		\rowcolor{headerblue}
		\headercell{Servicio} & \headercell{Casos} & \headercell{Qué verifican} \\
		\hline
		\endfirsthead
		\hline
		\rowcolor{headerblue}
		\headercell{Servicio} & \headercell{Casos} & \headercell{Qué verifican} \\
		\hline
		\endhead
		\texttt{quejas-service} & 112 & Validadores del registro público: correo en dos niveles, fechas,
		nombres, boleta, tutor y archivos por firma binaria; edición, corrección y cancelación; conciliación
		(RN-Q01 a RN-Q12) \\ \hline
		\texttt{primer-contacto-service} & 55 & Citas sin empalmes, plazo y rango máximo de consulta,
		token de respuesta; dictamen con Subdefensoría caída; sincronización de estatus; motor de
		antecedentes (RN-PC01 a RN-PC04) \\ \hline
		\texttt{admin-service} & 36 & Cuentas del personal y protección de la última cuenta de
		administración; bitácora con IP real; plantillas; respaldos en carpeta temporal; tablero (RN-A01 a
		RN-A04, RNF-A07) \\ \hline
		\texttt{auth-service} & 36 & Inicio de sesión, recuperación con código (cinco intentos, diez
		minutos), política de contraseñas, activación y emisión del token (RN-Q06 a RN-Q08) \\ \hline
		\texttt{revision-service} & 36 & Toma atómica entre recepcionistas, regreso, turnado idempotente,
		ventanilla e historial (RN-R01 a RN-R07, RNF-R06) \\ \hline
		\texttt{denunciado-service} & 29 & Todas las reglas del registro de la versión de la persona
		señalada (RF-D01) \\ \hline
		\texttt{subdefensoria-service} & 24 & Días hábiles (10 y 5), semáforo y alertas de vencimiento,
		archivo de respuesta (RN-SD01, RN-SD02) \\ \hline
		\texttt{notificaciones-service} & 20 & Centro de notificaciones e inyección de encabezados en el
		correo (CWE-93) \\ \hline
		\texttt{catalogo-service} & 8 & Alta, edición e importación de Excel tolerante a errores
		(RNF-A04) \\ \hline
		\texttt{chatbot-service} & 8 & Menú agrupado y administración de preguntas \\ \hline
		\textbf{Total} & \textbf{364} & \\ \hline
	\end{longtable}
	\normalsize

	\textbf{Resultado.} Los 364 casos se ejecutaron con Maven y resultaron aprobados (0 fallas, 0
	errores). La única prueba que requería base de datos ---la carga del contexto de Primer Contacto---
	se etiquetó como de integración y se excluye de la ejecución unitaria, para que la suite no dependa
	del entorno. Posteriormente se incorporaron ocho casos de resiliencia, descritos en la
	Sección~\ref{sec:pruebas_fiabilidad}; al cierre documentado se señala que
	FIA-05 y FIA-06 están implementadas, pero pendientes de su primera ejecución.
	Por esa razón, esos ocho casos no se suman a las 364 pruebas aprobadas.

	\subsection{Pruebas de integración y de sistema}
	La Tabla~\ref{tab:pruebas_integracion} lista los casos de integración y de sistema de prioridad P1 y
	P2 con su resultado. Los que tienen evidencia se ejecutaron durante el desarrollo de cada módulo con
	guiones y con el navegador; los marcados como pendientes forman parte de la suite de sistema que se
	ejecutará en el entorno de pruebas antes de la entrega a la Defensoría.

	\rowcolors{2}{rowblue}{white}
	\small
	\renewcommand{\arraystretch}{1.2}
	\setlength{\tabcolsep}{4pt}
	\begin{longtable}{|p{1.2cm}|L{6.2cm}|L{2.2cm}|L{2.0cm}|L{2.9cm}|}
		\hiderowcolors\caption{Pruebas de integración y de sistema}
		\label{tab:pruebas_integracion} \\ \showrowcolors
		\hline
		\rowcolor{headerblue}
		\headercell{ID} & \headercell{Qué se verifica} & \headercell{Requerimiento} &
		\headercell{Tipo} & \headercell{Resultado} \\
		\hline
		\endfirsthead
		\hline
		\rowcolor{headerblue}
		\headercell{ID} & \headercell{Qué se verifica} & \headercell{Requerimiento} &
		\headercell{Tipo} & \headercell{Resultado} \\
		\hline
		\endhead
		IT-06 & Dos recepcionistas abren la misma queja al mismo tiempo; solo una la obtiene &
		RNF-R06, RN-R03 & Dinámica, CN & Aprobada (10 de 10 intentos concurrentes) \\ \hline
		IT-20 & Activar una cuenta solo con un par folio--correo válido & RN-Q06 & Dinámica, CB &
		Aprobada en prueba unitaria; dinámica pendiente \\ \hline
		IT-22 & La queja turnada aparece en la bandeja de Primer Contacto & CU-REC-05, CU-PC-01 & Dinámica,
		CN & Aprobada en el sistema desplegado \\ \hline
		IT-23 & El dictamen competente crea el expediente en Subdefensoría & CU-PC-06, CU-SD-01 & Dinámica,
		CN & Aprobada en entorno local con la base de pruebas \\ \hline
		IT-24 & El cambio de estatus en Primer Contacto se refleja en la tabla de quejas & RN-R01 &
		Dinámica, CB & Aprobada en prueba unitaria de sincronización \\ \hline
		IT-09 & Los indicadores del panel de la Defensora coinciden con datos sembrados en PostgreSQL~16 &
		CU-ADM-03 & Dinámica, CB & Aprobada \\ \hline
		IT-11 & Inicio de sesión unificado: siete escenarios (sin sesión, enlace antiguo, Defensora en
		Recepción, cierre de sesión, recepcionista hacia el panel de la Defensora, entre otros) & CU-ADM-07,
		RNF-R02 & Dinámica, CN & Aprobada (7 de 7) \\ \hline
		IT-12 & Registro de la persona señalada con datos válidos, nombres con dígitos, ejecutable
		renombrado, credencial PDF o de 4~MB, tres credenciales y sin aviso & RF-D01, RNF-G08 & Dinámica,
		CB & Aprobada (19 de 19) \\ \hline
		REN-02 & Carga masiva de quejas sintéticas por el registro público, con credencial y evidencias &
		CU-QJ-07 & Dinámica, CN & Aprobada (170 quejas) \\ \hline
		ST-01 a ST-14 & Recorridos de extremo a extremo por actor, automatizados con Playwright & Casos de
		uso & Dinámica, CN & Pendiente: suite por automatizar; recorridos verificados manualmente \\ \hline
	\end{longtable}
	\normalsize

	\textbf{Compilación de producción.} Las cinco aplicaciones web se compilan con la configuración de
	producción, que falla si un componente excede su presupuesto de tamaño o si una plantilla hace
	referencia a un elemento inexistente. Las cinco compilan sin errores.

	\section{Pruebas de mantenibilidad}
	\label{sec:pruebas_mantenibilidad}
	La mantenibilidad se verificó con pruebas estáticas de caja blanca sobre el código Java de ocho de
	los diez servicios (225 archivos), y con la cobertura de las pruebas unitarias. Primer Contacto y
	Subdefensoría se analizaron con las mismas reglas de forma informativa.

	\subsection{Guía de estilo}
	El proyecto adoptó como guía de estilo el Google Java Style Guide \cite{google_java_style}, con cuatro
	desviaciones que reflejan lo que el código ya seguía de forma consistente, para no reformatear todo
	el proyecto sin aportar calidad ---sangría de 4 espacios en lugar de 2, líneas de 120 columnas en lugar de 100, guion bajo
	permitido en los nombres de paquete y el registrador \texttt{log} en minúsculas---. La guía se codificó como
	configuración de Checkstyle: las reglas \textbf{obligatorias} deben cumplirse sin excepción y las
	\textbf{recomendadas} se reportan como deuda técnica.



	La Tabla~\ref{tab:pruebas_mantenibilidad} resume los casos y sus resultados.

	\rowcolors{2}{rowblue}{white}
	\small
	\renewcommand{\arraystretch}{1.2}
	\setlength{\tabcolsep}{4pt}
	\begin{longtable}{|p{1.2cm}|L{5.5cm}|L{3.1cm}|L{4.7cm}|}
		\hiderowcolors\caption{Pruebas de mantenibilidad}
		\label{tab:pruebas_mantenibilidad} \\ \showrowcolors
		\hline
		\rowcolor{headerblue}
		\headercell{ID} & \headercell{Qué se verifica} & \headercell{Criterio} & \headercell{Resultado} \\
		\hline
		\endfirsthead
		\hline
		\rowcolor{headerblue}
		\headercell{ID} & \headercell{Qué se verifica} & \headercell{Criterio} & \headercell{Resultado} \\
		\hline
		\endhead
		MAN-01 & Reglas obligatorias de la guía: nombre de archivo, sin tabuladores, sin importaciones con
		comodín, llaves siempre, un enunciado por línea, una variable por declaración, nombres, sin
		\texttt{System.out} ni excepciones ignoradas, \texttt{equals} con \texttt{hashCode} & 0
		violaciones & Fallida: 4 violaciones en 225 archivos (tres en \texttt{RevisionQuejaService}: dos
		enunciados en una línea, dos variables en una declaración y un \texttt{if} sin llaves; una
		importación estática con comodín en \texttt{RespuestaDenunciadoService}) \\ \hline
		MAN-02 & Métricas: complejidad ciclomática $\leq 15$, NPath $\leq 400$, método $\leq 80$ líneas,
		$\leq 8$ parámetros, acoplamiento de salida $\leq 30$, archivo $\leq 800$ líneas & Se reporta como
		deuda & Informativa: 5 métodos con complejidad mayor a 15, 117 tipos públicos sin Javadoc, 90 líneas
		de más de 120 columnas \\ \hline
		MAN-03 & Cobertura de líneas en los paquetes de servicios, validaciones y utilerías, medida con
		JaCoCo & $\geq 70~\%$ por servicio & Pendiente de su primera ejecución (el reporte se
		genera en cada ejecución de las pruebas) \\ \hline
		MAN-04 & Cada servicio publica su contrato OpenAPI & 10 de 10 & Aprobada (Primer Contacto y
		Subdefensoría lo incorporaron al cierre) \\ \hline
	\end{longtable}
	\normalsize

	El cumplimiento global de las reglas obligatorias es de cuatro violaciones en 225 archivos: corregirlas
	permite aprobar MAN-01. La deuda reportada por MAN-02 no es bloqueante, pero indica dónde conviene
	concentrar la refactorización; los métodos más complejos están en los servicios de revisión y de
	seguimiento, que concentran la lógica de los tableros.

	\section{Pruebas de usabilidad}
	\label{sec:pruebas_usabilidad}
	La usabilidad no puede verificarse por completo con herramientas automáticas, porque depende de la
	experiencia de personas concretas. Se combinaron por ello mediciones automáticas de accesibilidad con
	pruebas con usuarios. La Tabla~\ref{tab:pruebas_usabilidad} resume los casos.

	\rowcolors{2}{rowblue}{white}
	\small
	\renewcommand{\arraystretch}{1.2}
	\setlength{\tabcolsep}{4pt}
	\begin{longtable}{|p{1.2cm}|L{5.6cm}|L{2.2cm}|L{2.6cm}|L{2.9cm}|}
		\hiderowcolors\caption{Pruebas de usabilidad}
		\label{tab:pruebas_usabilidad} \\ \showrowcolors
		\hline
		\rowcolor{headerblue}
		\headercell{ID} & \headercell{Qué se verifica} & \headercell{Requerimiento} &
		\headercell{Criterio} & \headercell{Resultado} \\
		\hline
		\endfirsthead
		\hline
		\rowcolor{headerblue}
		\headercell{ID} & \headercell{Qué se verifica} & \headercell{Requerimiento} &
		\headercell{Criterio} & \headercell{Resultado} \\
		\hline
		\endhead
		USA-01 & Cuestionario SUS a cinco usuarios después de las pruebas alfa \cite{brooke_sus} & RNF-G09 &
		Puntaje $\geq 68$ & Pendiente: sesión alfa con la Defensoría \\ \hline
		USA-02 & Tareas cronometradas: presentar una queja, turnar, agendar una cita y dictaminar & RNF-Q05,
		RNF-A05 & Tiempo y errores por tarea & Pendiente: sesión alfa \\ \hline
		USA-03 & Accesibilidad básica de los portales: contraste, etiquetas y uso con teclado, medida con
		Lighthouse & RNF-G09 & Puntuación $\geq 0.85$ & Configurada; informativa hasta definir el umbral
		por portal \\ \hline
		USA-04 & El portal del quejoso permite presentar la queja completa desde un teléfono & RNF-Q01 &
		Flujo completo & Pendiente de verificación formal; diseño adaptable implementado \\ \hline
		USA-05 & El formulario de queja se presenta en pasos numerados con previsualización de archivos &
		RNF-Q05 & Cumple & Aprobada por inspección \\ \hline
		USA-06 & Toda operación muestra éxito o error; sin fallos silenciosos & RNF-Q07 & Cumple & Aprobada
		por inspección; se propone automatizarla con Playwright \\ \hline
		USA-07 & Cada indicador del panel de la Defensora abre la lista filtrada & RNF-A05 & Cumple &
		Aprobada en revisión con capturas de las cinco pantallas \\ \hline
	\end{longtable}
	\normalsize

	La validación cualitativa más importante ocurrió durante el desarrollo: las sesiones de revisión con
	el personal produjeron cambios concretos ---la redacción jurídica de los requisitos de admisión, el
	paso de dos a cinco preguntas en la validación, la matrícula en dos pasos y el rediseño completo de
	Primer Contacto y del panel de la Defensora---, lo que equivale a varias iteraciones de prueba de
	usabilidad informal. La medición cuantitativa con SUS queda como criterio de salida de la etapa de
	aceptación.

	\section{Pruebas de fiabilidad}
	\label{sec:pruebas_fiabilidad}
	Las pruebas de fiabilidad verifican la tolerancia a fallos ---que el fallo de un componente no
	detenga a los demás--- y la recuperabilidad ---que los datos puedan restaurarse---. Las primeras se
	implementaron como pruebas unitarias que simulan un servicio caído; las segundas, como guiones que
	operan sobre bases desechables.

	\rowcolors{2}{rowblue}{white}
	\small
	\renewcommand{\arraystretch}{1.2}
	\setlength{\tabcolsep}{4pt}
	\begin{longtable}{|p{1.2cm}|L{5.6cm}|L{2.2cm}|L{2.6cm}|L{2.9cm}|}
		\hiderowcolors\caption{Pruebas de fiabilidad}
		\label{tab:pruebas_fiabilidad} \\ \showrowcolors
		\hline
		\rowcolor{headerblue}
		\headercell{ID} & \headercell{Qué se verifica} & \headercell{Requerimiento} &
		\headercell{Criterio} & \headercell{Resultado} \\
		\hline
		\endfirsthead
		\hline
		\rowcolor{headerblue}
		\headercell{ID} & \headercell{Qué se verifica} & \headercell{Requerimiento} &
		\headercell{Criterio} & \headercell{Resultado} \\
		\hline
		\endhead
		FIA-05 & Con el servicio de notificaciones sin conexión o respondiendo 500, los avisos de queja
		creada, al tutor, de cambio de estatus y al personal no propagan la excepción (4 casos) & RNF-Q04 &
		Ninguna excepción se propaga & Implementada; pendiente de su primera ejecución \\ \hline
		FIA-06 & Con el servicio de notificaciones caído, el correo de observaciones y los avisos de
		Recepción no impiden el regreso de la queja (4 casos) & RNF-R04 & Ninguna excepción se propaga &
		Implementada; pendiente de su primera ejecución \\ \hline
		FIA-07 & La importación del catálogo con filas mal formadas reporta un resumen y no falla en bloque
		& RNF-A04 & Resumen de creadas, actualizadas y con error & Aprobada \\ \hline
		FIA-04 & Con Subdefensoría caída al dictaminar, el dictamen se conserva y puede reenviarse &
		RN-PC03 & El dictamen no se pierde & Aprobada (prueba unitaria) \\ \hline
		FIA-03 & Respaldo manual, descarga y restauración en una base desechable & RNF-G04, CU-ADM-05 & Sin
		errores de \texttt{psql}; tablas principales con filas & Pendiente de ejecución en el entorno de
		pruebas \\ \hline
		FIA-01 & Detener un contenedor auxiliar mientras se usa el sistema & RNF-Q04 & El resto opera &
		Aprobada en operación (incidentes de la Sección~\ref{sec:incidentes}) \\ \hline
		FIA-02 & Reiniciar el servidor de aplicación & RNF-G04 & Los contenedores se levantan solos &
		Aprobada en operación \\ \hline
	\end{longtable}
	\normalsize

	Las pruebas FIA-05 y FIA-06 son el ejemplo más claro de cómo un incidente se convirtió en prueba. El
	defecto descrito en la Sección~\ref{sec:incidentes} ---peticiones colgadas por una llamada sin tiempo
	de espera--- se habría detectado con una prueba que simulara el servicio caído. Las pruebas usan un
	cliente HTTP simulado que lanza un error de conexión o un error 500 y verifican que ningún método del
	cliente de notificaciones propague la excepción hacia el servicio que lo invoca. La restauración
	(FIA-03) es la prueba de mayor impacto de esta sección: un respaldo que no puede restaurarse no es un
	respaldo, y la única forma de saberlo es restaurarlo. El guion genera un respaldo desde el panel,
	lo descarga, lo restaura en una base nueva y verifica que las tablas de quejas, personal y
	dependencias tengan filas.

	\section{Pruebas de seguridad}
	\label{sec:pruebas_seguridad}
	La seguridad se verificó en dos planos complementarios. En el plano estático, se extrajo del código la
	matriz completa de roles por punto de acceso y se comparó con las reglas que exigen los
	requerimientos. En el plano dinámico, la misma matriz se recorre contra el sistema en ejecución con un
	token por rol.

	\subsection{Matriz de permisos extraída del código}
	Un guion recorre los controladores de ocho de los diez servicios ---Primer Contacto y Subdefensoría
	declaran sus permisos solo en la cadena de filtros y se verifican en la prueba dinámica---, lee las rutas, los
	métodos HTTP, las anotaciones \texttt{@PreAuthorize} y las reglas públicas de la cadena de filtros, y
	genera una matriz de 92 puntos de acceso con los roles que cada uno admite. Un segundo guion compara
	esa matriz con un archivo de reglas que traduce los requerimientos a expresiones sobre las rutas ---por
	ejemplo, ``\texttt{/api/admin/seguridad/(respaldos|restaurar)} solo admite a la Defensora''--- y
	reporta cada diferencia. Al ser estática, la prueba puede ejecutarse con cada cambio del código y detecta, antes del
	despliegue, un punto de acceso nuevo que se haya publicado con permisos de más.

	La primera ejecución encontró seis diferencias. Cuatro correspondían a los respaldos y la
	restauración, que el código abre al rol \texttt{DEFENSOR} mientras que la redacción original del
	requerimiento RNF-A06 los reservaba a un rol técnico de administración. Las otras dos
	correspondían a la consulta y emisión de acuerdos de conciliación, que el código abre también a
	Subdefensoría y a la Defensora, mientras que la redacción original de RN-R02 reservaba todo el
	servicio de revisión a Recepción. En ambos casos el análisis concluyó que el código reflejaba la
	decisión de diseño vigente ---por convención con la Defensoría, la administración corresponde a la
	Defensora, y la conciliación es una actuación de varias áreas---, de modo que se corrigió la redacción de los requerimientos
	(Tablas~\ref{tab:rnf} y~\ref{tab:reglas_negocio}) y el archivo de reglas, y no el código. El
	resultado ilustra un beneficio de esta prueba que va más allá de la seguridad: obliga a que el
	documento de requerimientos y el código digan lo mismo.

	\subsection{Casos de prueba de seguridad}
	La Tabla~\ref{tab:pruebas_seguridad} reúne los casos de seguridad con su resultado.

	\rowcolors{2}{rowblue}{white}
	\small
	\renewcommand{\arraystretch}{1.2}
	\setlength{\tabcolsep}{4pt}
	\begin{longtable}{|p{1.2cm}|L{5.6cm}|L{2.2cm}|L{2.6cm}|L{2.9cm}|}
		\hiderowcolors\caption{Pruebas de seguridad}
		\label{tab:pruebas_seguridad} \\ \showrowcolors
		\hline
		\rowcolor{headerblue}
		\headercell{ID} & \headercell{Qué se verifica} & \headercell{Requerimiento} &
		\headercell{Criterio} & \headercell{Resultado} \\
		\hline
		\endfirsthead
		\hline
		\rowcolor{headerblue}
		\headercell{ID} & \headercell{Qué se verifica} & \headercell{Requerimiento} &
		\headercell{Criterio} & \headercell{Resultado} \\
		\hline
		\endhead
		SEG-07 & Matriz de permisos extraída del código contra las reglas de los requerimientos (34 puntos
		de acceso regulados de 92) & RNF-A06, RN-R02, RNF-R02, RNF-Q01, RNF-Q02, RN-Q12 & 0 diferencias &
		Fallida en la línea base (6 diferencias); resuelta al alinear los requerimientos \\ \hline
		SEG-01 & La matriz en vivo: sin token, 401 o 403; con un rol no permitido, 403; con el permitido,
		la petición pasa el filtro & RN-R02, RNF-A06 & 0 accesos indebidos & Pendiente de ejecución en el
		entorno de pruebas \\ \hline
		SEG-10 & Las rutas de Primer Contacto y Subdefensoría exigen token (antes eran públicas) & RN-R02 &
		403 sin token & Aprobada en el servidor: consulta, PDF de oficios y evidencias responden 403 sin
		token \\ \hline
		IT-04 & Un quejoso no puede consultar ni editar la queja de otra persona cambiando el folio & RN-Q11
		& 404 & Aprobada por revisión de código (búsqueda por folio y correo del token) \\ \hline
		SEG-08 & El token vence a la hora (\texttt{exp} $-$ \texttt{iat} $= 3\,600$) y un token alterado se
		rechaza & RNF-Q06 & 3\,600~s; 401 o 403 & Aprobada por configuración y prueba unitaria de emisión;
		dinámica pendiente \\ \hline
		SEG-09 & Restaurar sin \texttt{"confirmar": true} se rechaza & RNF-A03 & Respuesta 4xx & Probable falla
		por revisión de código: se rechaza con una excepción genérica que se traduce en 500, no en 400 \\ \hline
		SEG-03 & Recorrido de rutas (\texttt{../}) en la descarga y restauración de respaldos & RNF-A06 &
		Distinto de 200 & Fallida por revisión de código: el nombre de archivo no se valida contra la carpeta
		de respaldos (riesgo bajo: solo cuentas de administración) \\ \hline
		SEG-05 & Secretos ---contraseña de la base, clave de firma de tokens, credenciales de correo--- en el
		repositorio & RNF-G07 & 0 hallazgos & Fallida: residen en texto plano en los archivos de
		configuración \\ \hline
		SEG-04 & Dependencias Maven y npm con vulnerabilidades altas o críticas con parche & RNF-G07 & 0 &
		Pendiente de su primera ejecución \\ \hline
		SEG-02 & Escaneo pasivo de vulnerabilidades web de los cinco portales & RNF-G06 & Sin alertas
		altas & Pendiente (informativa) \\ \hline
		SEG-11 & Los puertos de los servicios y de la base no son alcanzables desde Internet & RNF-G06 &
		Conexión rechazada & Aprobada en operación (reglas en la tabla \texttt{raw}) \\ \hline
	\end{longtable}
	\normalsize

	Tres resultados requieren acción. El más importante es SEG-05: las contraseñas y la clave de firma
	de los tokens están en los archivos de configuración versionados, y la rama de producción del
	repositorio pudo clonarse sin credenciales. La corrección consiste en rotar esos secretos, pasarlos a
	variables de entorno del servidor y retirarlos del historial del repositorio. SEG-03 se corrige
	verificando que la ruta normalizada del archivo permanezca dentro de la carpeta de respaldos, y SEG-09
	con un manejador que traduzca la falta de confirmación a una respuesta 400.

	\section{Pruebas de capacidad}
	\label{sec:pruebas_capacidad}
	La capacidad se verifica en dos dimensiones: el tamaño máximo de una petición, que determina cuántas
	evidencias puede adjuntar un quejoso, y el número de usuarios concurrentes. Las pruebas de tamaño
	envían formularios multiparte vacíos a propósito, con un relleno del tamaño deseado, de modo que no
	crean quejas en el entorno.

	\rowcolors{2}{rowblue}{white}
	\small
	\renewcommand{\arraystretch}{1.2}
	\setlength{\tabcolsep}{4pt}
	\begin{longtable}{|p{1.2cm}|L{5.6cm}|L{2.2cm}|L{2.6cm}|L{2.9cm}|}
		\hiderowcolors\caption{Pruebas de capacidad}
		\label{tab:pruebas_capacidad} \\ \showrowcolors
		\hline
		\rowcolor{headerblue}
		\headercell{ID} & \headercell{Qué se verifica} & \headercell{Requerimiento} &
		\headercell{Criterio} & \headercell{Resultado} \\
		\hline
		\endfirsthead
		\hline
		\rowcolor{headerblue}
		\headercell{ID} & \headercell{Qué se verifica} & \headercell{Requerimiento} &
		\headercell{Criterio} & \headercell{Resultado} \\
		\hline
		\endhead
		CAP-01 & Una petición de 95~MB atraviesa el proxy & RNF-Q03 & Distinto de 413 & Pendiente de
		ejecución \\ \hline
		CAP-02 & Una petición de 105~MB es rechazada por el proxy & RNF-Q03 & 413 & Pendiente de ejecución \\ \hline
		CAP-03, CAP-04 & Lo mismo directamente contra el servicio de quejas & RNF-Q03 & Distinto de 413 / 413
		& Pendiente de ejecución \\ \hline
		CAP-06 & Coherencia entre el requerimiento y la configuración de cada servicio & RNF-Q03 & Iguales
		& Observada y resuelta: el requerimiento indicaba 60~MB y los servicios aceptan 100~MB; se actualizó
		el requerimiento \\ \hline
		CAP-07 & Una queja con seis archivos, uno de 12.5~MB, se registra sin terminar el contenedor &
		RNF-Q03 & Respuesta 200 & Aprobada tras el redimensionamiento de memoria \\ \hline
		CAP-05 & 50 usuarios virtuales concurrentes sobre las consultas públicas (catálogo, menú del
		chatbot, consulta por folio) & RNF-G05 & p95 $< 1$~s; errores $< 1~\%$ & Pendiente de ejecución
		con k6 \\ \hline
	\end{longtable}
	\normalsize

	La prueba CAP-07 merece un comentario. El límite de tamaño de una petición no es solo una cifra de
	configuración: como las evidencias se guardan como contenido binario en la base, cada megabyte
	recibido se copia varias veces en memoria. Un límite de 100~MB por petición exige que el servicio
	tenga memoria para varias copias de esa petición, y la prueba con archivos reales fue la que reveló
	que el contenedor original no la tenía. Por eso RNF-Q03 se verifica con dos tipos de prueba: la de
	límites, que comprueba el rechazo, y la de archivos reales, que comprueba que lo aceptado se procesa.

	\section{Pruebas de confidencialidad}
	\label{sec:pruebas_confidencialidad}
	La confidencialidad se verifica con pruebas de caja negra que observan, desde fuera del servidor, cómo
	viaja la información. Son pasivas ---no modifican datos--- y por ello pueden ejecutarse también contra
	producción. El guion usa \texttt{openssl} para negociar conexiones con cada versión de TLS y
	\texttt{curl} para revisar redirecciones y encabezados.

	\rowcolors{2}{rowblue}{white}
	\small
	\renewcommand{\arraystretch}{1.2}
	\setlength{\tabcolsep}{4pt}
	\begin{longtable}{|p{1.2cm}|L{5.6cm}|L{2.2cm}|L{2.6cm}|L{2.9cm}|}
		\hiderowcolors\caption{Pruebas de confidencialidad}
		\label{tab:pruebas_confidencialidad} \\ \showrowcolors
		\hline
		\rowcolor{headerblue}
		\headercell{ID} & \headercell{Qué se verifica} & \headercell{Requerimiento} &
		\headercell{Criterio} & \headercell{Resultado} \\
		\hline
		\endfirsthead
		\hline
		\rowcolor{headerblue}
		\headercell{ID} & \headercell{Qué se verifica} & \headercell{Requerimiento} &
		\headercell{Criterio} & \headercell{Resultado} \\
		\hline
		\endhead
		CON-01 & Una petición HTTP se redirige a HTTPS & RNF-Q02 & 301 a \texttt{https} & Aprobada por
		configuración; dinámica pendiente \\ \hline
		CON-02 & Se aceptan TLS 1.2 y 1.3 & RNF-Q02 & Negociación exitosa & Aprobada por configuración
		(\texttt{ssl\_protocols}) \\ \hline
		CON-03 & Se rechazan TLS 1.0 y 1.1 & RNF-Q02 & Negociación fallida & Aprobada por configuración \\ \hline
		CON-04 & Certificado válido con al menos 14 días de vigencia & RNF-Q02 & Cumple & Aprobada por
		configuración (renovación automática con Certbot) \\ \hline
		CON-05 & El servidor envía el encabezado \texttt{Strict-Transport-Security} & RNF-Q02 & Presente &
		Fallida: el proxy no lo envía \\ \hline
		CON-06 & Las respuestas JSON de la API no se comprimen (mitigación de BREACH) & RNF-Q02 & Sin
		\texttt{Content-Encoding} & Aprobada por configuración \\ \hline
		CON-07 & Las rutas internas de notificaciones e ingesta responden 403 desde Internet & RNF-G06 &
		403 & Aprobada por configuración; dinámica pendiente \\ \hline
		CON-08 & El inicio de sesión no usa \textit{cookies} ni acepta credenciales en la URL; la sesión
		viaja en el encabezado \texttt{Authorization} & RNF-Q02 & Sin \texttt{Set-Cookie} & Aprobada por
		diseño \\ \hline
		CON-09 & Las evidencias de una queja solo se entregan al quejoso dueño o al personal autorizado &
		RN-Q11 & 403 o 404 para terceros & Aprobada (cubierta por IT-04 y SEG-10) \\ \hline
	\end{longtable}
	\normalsize

	La única falla, CON-05, se corrige con una línea en la configuración del servidor seguro del proxy.
	Varias pruebas figuran como aprobadas por configuración porque la configuración del proxy se revisó
	línea por línea; su ejecución dinámica, que confirma que esa configuración es la que realmente está
	cargada, queda pendiente para el entorno de pruebas.
	
	\section{Pruebas de los componentes de inteligencia artificial}
\label{sec:pruebas_ia}

La verificación de los componentes de inteligencia artificial comprendió dos tipos de actividades:
la evaluación experimental de la recuperación de antecedentes con datos sintéticos y las
pruebas funcionales y de integración realizadas sobre los servicios desplegados. En el caso de
los resúmenes se verificaron tanto la respuesta técnica del servicio como la presencia de
información relevante en una salida generada. Se distinguen expresamente las mediciones
experimentales de las comprobaciones puntuales de funcionamiento, pues una respuesta HTTP
correcta no demuestra por sí sola la calidad del contenido producido.

\subsection{Evaluación de recuperación de antecedentes}
\label{sec:pruebas_ia_antecedentes}

Se utilizaron dos conjuntos sintéticos. El primero estuvo compuesto por 200 quejas extensas,
organizadas en 25 escenarios con ocho variantes de redacción por escenario; cada queja se
empleó como consulta y la relevancia se determinó mediante su pertenencia al mismo escenario.
Para interpretar correctamente esta evaluación, debe comprobarse y dejarse explícito
si cada consulta fue excluida de la colección de candidatos. Si no se excluyó,
una coincidencia consigo misma podría sobrestimar la recuperación. La información
presentada no basta para establecer cuál fue el tratamiento aplicado.
El segundo incluyó 15 consultas contrastivas frente a 45 documentos candidatos, con un
antecedente positivo sintético por consulta y documentos distractores. Estas etiquetas
representan relaciones diseñadas para el experimento y no equivalen a una valoración jurídica
de los expedientes.

Para estudiar la combinación de los recuperadores léxico (TF-IDF) y semántico (MiniLM) se
reutilizaron los rankings calculados previamente por ambos modelos. La fusión se realizó con
\textit{Reciprocal Rank Fusion} (RRF), considerando los diez primeros candidatos de cada
recuperador y una constante de suavizado de 60. Se examinaron 101 distribuciones de peso,
desde 100\,\% léxico hasta 100\,\% semántico, con incrementos de un punto porcentual. Esta
operación recalculó la fusión sobre rankings guardados; no supuso ejecutar nuevamente MiniLM
para cada combinación ni consultar los expedientes del servidor de producción.

Se calcularon precisión en los primeros cinco resultados (P@5), exhaustividad en los primeros
diez (R@10), rango recíproco medio (MRR) y ganancia acumulada descontada normalizada
(nDCG@5), entre otras métricas. La Tabla~\ref{tab:ia_pesos_200} muestra una selección de
ponderaciones para las 200 quejas extensas.

\begin{longtable}{|p{3.8cm}|p{2.6cm}|p{2.6cm}|p{2.6cm}|}
\hiderowcolors\caption{Resultados de recuperación en 200 quejas sintéticas extensas}
\label{tab:ia_pesos_200} \\ \showrowcolors
\hline
\textbf{TF-IDF / MiniLM} & \textbf{P@5} & \textbf{R@10} & \textbf{nDCG@5} \\
\hline
\endfirsthead
\hline
\textbf{TF-IDF / MiniLM} & \textbf{P@5} & \textbf{R@10} & \textbf{nDCG@5} \\
\hline
\endhead
100\,\% / 0\,\% & 45.2\,\% & 91.5\,\% & 0.435 \\ \hline
90\,\% / 10\,\% & 52.3\,\% & 91.5\,\% & 0.496 \\ \hline
75\,\% / 25\,\% & 52.0\,\% & 91.5\,\% & 0.490 \\ \hline
50\,\% / 50\,\% & 49.7\,\% & 67.5\,\% & 0.471 \\ \hline
25\,\% / 75\,\% & 48.5\,\% & 46.5\,\% & 0.457 \\ \hline
0\,\% / 100\,\% & 33.5\,\% & 46.5\,\% & 0.325 \\ \hline
\end{longtable}

La búsqueda con predominio léxico obtuvo el mejor P@5 entre las configuraciones mostradas,
pero la diferencia entre 90\,/\,10 y 75\,/\,25 fue de solo 0.3 puntos porcentuales. En la
búsqueda más fina de 101 combinaciones, la ponderación 91\,/\,9 también alcanzó un P@5
de 52.3\,\%. Por tanto, los datos no permiten afirmar una superioridad concluyente de una
ponderación concreta. Además, el resultado depende de la métrica seleccionada: el orden más
favorable para P@5 no tiene por qué serlo para la primera posición.

La configuración final adoptada para el servicio fue 75\,\% TF-IDF y
25\,\% MiniLM en la fusión RRF ponderada. En este primer conjunto alcanzó
P@5 de 52.0\,\%, R@10 de 91.5\,\% y nDCG@5 de 0.490. La elección
no equivale a afirmar que sea la mejor ponderación para todos los escenarios: se
conserva como configuración de implementación y se reportan sus limitaciones.

La Tabla~\ref{tab:ia_pesos_contraste} resume la prueba contrastiva. En este conjunto, donde
se introdujeron diferencias de redacción y distractores, la mayor influencia del recuperador
semántico produjo una exhaustividad sustancialmente superior. Dado que solo existe
un antecedente relevante por consulta, el valor máximo teórico de P@5 es 20\,\%
cuando se devuelven cinco resultados: por ello, sus porcentajes no son directamente
comparables con los del conjunto de 200 quejas, que tiene varios documentos
relevantes por escenario. Para este segundo conjunto resultan particularmente
interpretables las métricas Hit@k y MRR; deben añadirse sus resultados solo
cuando se cuente con las salidas verificadas del experimento.

\begin{longtable}{|p{3.8cm}|p{2.6cm}|p{2.6cm}|p{2.6cm}|}
\hiderowcolors\caption{Resultados en 15 consultas sintéticas contrastivas}
\label{tab:ia_pesos_contraste} \\ \showrowcolors
\hline
\textbf{TF-IDF / MiniLM} & \textbf{P@5} & \textbf{R@10} & \textbf{nDCG@5} \\
\hline
\endfirsthead
\hline
\textbf{TF-IDF / MiniLM} & \textbf{P@5} & \textbf{R@10} & \textbf{nDCG@5} \\
\hline
\endhead
90\,\% / 10\,\% & 2.7\,\% & 13.3\,\% & 0.084 \\ \hline
75\,\% / 25\,\% & 2.7\,\% & 13.3\,\% & 0.084 \\ \hline
50\,\% / 50\,\% & 9.3\,\% & 86.7\,\% & 0.216 \\ \hline
25\,\% / 75\,\% & 9.3\,\% & 93.3\,\% & 0.219 \\ \hline
0\,\% / 100\,\% & 17.3\,\% & 93.3\,\% & 0.669 \\ \hline
\end{longtable}

En particular, la ponderación final 75/25 alcanzó P@5 de 2.7\,\%,
R@10 de 13.3\,\% y nDCG@5 de 0.084 en las consultas contrastivas.
Este resultado advierte que la configuración seleccionada puede omitir
antecedentes relevantes cuando existe poca coincidencia léxica, por lo que
debe mantenerse la revisión humana y ampliarse la evaluación de casos
parafraseados.

Los resultados evidencian un compromiso entre coincidencia léxica y recuperación semántica,
no una ponderación universalmente óptima. La diferencia entre conjuntos podría
relacionarse con el grado de coincidencia de vocabulario, las paráfrasis y la
construcción de distractores, pero estas hipótesis requieren un análisis
individual de consultas fallidas para comprobarse. La selección de pesos se examinó también mediante
particiones por escenario, con el propósito de no ajustar y medir una configuración sobre los
mismos grupos sintéticos. No obstante, el documento no especifica el número de
particiones, sus tamaños ni los resultados diferenciados de selección y evaluación,
por lo que aún no puede considerarse reproducible dicho procedimiento. Aun así, los conjuntos no han sido etiquetados por especialistas
de la Defensoría; por tanto, estas medidas describen el comportamiento experimental del
recuperador y no su precisión sobre expedientes institucionales reales.

\subsection{Funcionamiento del índice y respuesta de antecedentes}
\label{sec:pruebas_ia_indice}

En el servidor se verificó la construcción del índice sobre 1\,210 registros y la
persistencia de sus representaciones semánticas. Una actualización posterior, sin cambios
en los documentos, informó 1\,210 representaciones reutilizadas y ninguna recalculada.
El servicio informó estado listo mediante su punto de comprobación de disponibilidad.

También se ejecutaron consultas HTTP de búsqueda de antecedentes con respuesta satisfactoria
(código 200). En una consulta sintética de revisión de calificación, la API de análisis
respondió en aproximadamente 0.135 segundos; en otra consulta realizada mediante el contrato
de integración se registró un tiempo cercano a 0.111 segundos. Estas cifras son observaciones
puntuales con el índice disponible, no percentiles de latencia ni pruebas de concurrencia.

\subsection{Generación de resúmenes e integración con los servicios}
\label{sec:pruebas_ia_resumenes}

Se envió una queja sintética al servicio de resúmenes, que realiza una selección extractiva
de información y posteriormente utiliza un modelo generativo para redactar un borrador en
estilo institucional. La solicitud devolvió HTTP 200 y un tiempo total de 26.304 segundos.
Se confirmó la presencia de un texto generado y del indicador de revisión humana requerida.

La inspección del resultado permitió identificar tanto elementos conservados como omisiones.
Es una observación cualitativa de un caso, no una tasa de fidelidad o de cobertura
para el componente de generación.
Se mantuvieron las referencias a los correos electrónicos y capturas del sistema, así como
la negación de agresión física; en cambio, no se conservaron las dos fechas explícitas ni
el cambio de calificación de ocho a seis. Por tanto, esta prueba fue satisfactoria en
respuesta técnica, pero detectó una limitación de cobertura factual. No se completó en esa
prueba la atribución de las omisiones a la fase extractiva o a la edición generativa.
Para efectuar esa atribución podría compararse, incluso con entradas sintéticas,
el texto original, la selección extractiva y el borrador generado, verificando
acontecimientos, negaciones, fechas, cantidades y evidencias. Hasta disponer de
esa comparación, no corresponde atribuir la pérdida de información exclusivamente
a MMR o a Qwen3.

La evaluación efectuada no incluyó un conjunto amplio de resúmenes evaluados
mediante criterios institucionales, ni una validación formal con expedientes
reales y personal de la Defensoría de los Derechos Politécnicos. Estas actividades
no pudieron completarse debido al tiempo disponible para el desarrollo y la
presentación del Trabajo Terminal. En consecuencia, las pruebas con narrativas
sintéticas permiten comprobar la integración y detectar errores puntuales, pero
no establecer un nivel de fidelidad, cobertura o utilidad generalizable a la
operación de la Defensoría. Como trabajo futuro, se propone una evaluación con
expedientes reales debidamente autorizados y revisión de personal institucional,
con criterios explícitos para la conservación de hechos, actores, fechas,
negaciones, cantidades, evidencias y solicitudes. Este trabajo debe distinguir
las posibles omisiones de la etapa extractiva de las modificaciones introducidas
por la etapa generativa.

Después del despliegue se comprobó, mediante solicitudes al punto de salud, que el servicio
Java de Primer Contacto podía comunicarse con el servicio de resúmenes (HTTP 200) y que
este último podía comunicarse con el servidor del modelo generativo (HTTP 200). Estas
comprobaciones acreditan conectividad entre componentes, pero no constituyen una medición
de calidad del resumen ni una prueba de carga.

\subsection{Alcance de los resultados de inteligencia artificial}

Las pruebas ejecutadas permiten documentar el comportamiento de la fusión de rankings sobre
conjuntos sintéticos, la reutilización del índice y el funcionamiento de las interfaces HTTP.
Sin embargo, no demuestran que la ponderación 75/25 sea óptima sobre quejas reales ni
la conservación factual sistemática de los resúmenes. La verificación de una
salida generada no sustituye un estudio con múltiples documentos y criterios de
revisión previamente definidos. Por cuestiones de tiempo no fue posible llevar
a cabo la validación con expedientes reales ni la revisión formal por parte del
personal de la Defensoría. Este alcance debe distinguirse de las pruebas técnicas
que sí fueron ejecutadas con datos sintéticos. La evaluación institucional queda
identificada como trabajo futuro y requiere acceso autorizado a los expedientes,
definición de criterios con los especialistas y análisis independiente de las
etapas extractiva y generativa. Tampoco se documentan una prueba de concurrencia
para la inferencia ni una distribución representativa de los tiempos de respuesta.

